\chapter{Textos de referencia del conjunto de evaluación}\label{anexo:textos_evaluacion}

Este apéndice recoge el texto de referencia completo de los 10 casos de prueba que integran el
conjunto de evaluación descrito en la Sección~\ref{section:conjunto_datos_evaluacion}, empleados
como dictado original en la evaluación \englishText{end-to-end} tanto en modalidad \textit{batch}
como \textit{streaming} (Sección~\ref{section:evaluacion_e2e}), así como una muestra de los
informes clínicos reales generados por el sistema a partir de estos textos
(Sección~\ref{anexo:ejemplos_informes_generados}).

\section{Caso 1}\label{anexo:texto_caso_1}

Paciente mujer de 53 años de edad, diagnosticada de hernia discal, en tratamiento con Durogesic
50 $\mu$g/h durante 6 años, con buena eficacia y seguridad. Se cambia por fentanilo genérico, y se nota
un empeoramiento del cuadro doloroso. Se vuelve a cambiar a Durogesic®, y de nuevo se estabiliza la
analgesia.

\section{Caso 2}\label{anexo:texto_caso_2}

Varón de 29 años. Hijo mediano de tres hermanos (los otros dos, están sanos). Casado y sin hijos.
No bebedor, ni fumador. Amigdalitis de repetición en la infancia, siendo operado a los 6 años.
Asintomático. Hábito intestinal normal. Analítica normal con ligera elevación persistente en las
transaminasas. TTG positiva (2,46). DQ2 (-) y DQ8 (+). La biopsia duodenal mostró atrofia moderada
de las vellosidades intestinales, acompañada de infiltrado inflamatorio importante, a nivel de la
lámina propia. Está con DSG desde hace 2 años, con buena respuesta clínica y analítica.

\section{Caso 3}\label{anexo:texto_caso_3}

Mujer de 42 años de edad con antecedentes personales de hernia umbilical intervenida, HTA.
Presentaba dolor en FII de 2 años de evolución, tenesmo rectal y heces acintadas. A la exploración
física destacaba dolor difuso a la palpación profunda en FII. En la analítica destacaba 16.000
leucocitos y 91\% neutrofilos. Se realizó enema opaco y colonoscopia en ambas destacaba una imagen
estenosante a nivel de sigma. El estudio histológico de las biopsias no fue concluyente. Se realizó
la resección de la zona estenosante con anastomosis término-terminal. El estudio histológico
infirmó la lesión de endometriosis de colon con afectación transmural ocasionando áreas de fibrosis
y de hemorragia.

\section{Caso 4}\label{anexo:texto_caso_4}

Paciente de 38 años diagnosticado de pancolitis ulcerosa desde el año 2004, corticodependiente. Ese
mismo año presenta un cuadro de actividad inflamatoria intestinal con diarrea sin sangre,
afectación importante del estado general y pérdida de peso de 15 kg en un mes, así como
intolerancia digestiva a tratamiento con azatioprina. Durante estos 3 últimos años ha precisado
varios ingresos debidos a graves brotes de diarrea (5-6 deposiciones diarias) con sangre roja.

En junio de 2007 presentó un brote prolongado de colitis ulcerosa con activación moderada y mala
respuesta al tratamiento farmacológico con criterios de cortirrefractariedad, momento en el que se
inicia la citoaféresis, recibiendo una sesión cada 7 días durante 5 semanas. Las vías de acceso
fueron las venas antecubitales de los brazos, y durante las sesiones no se presentaron
complicaciones ni reacciones adversas. Desde la segunda sesión presentó una notable mejoría clínica
y en la actualidad recibe tratamiento oral de micofenolato-mofetil.

\section{Caso 5}\label{anexo:texto_caso_5}

Mujer de 35 años acude a consulta por mala agudeza visual desde la infancia, sin tratamiento. A la
exploración se encontró cabello amarillento, al igual que cejas y pestañas; piel rojiza, sin nevos;
agudeza visual con corrección de 20/400 en ojo derecho 20/200 en ojo izquierdo; exotropía de 50
dioptrías, nistagmo horizontal; ambos ojos con presión intraocular de 12 mmHg, conjuntiva normal,
córnea con algunos depósitos estromales blanquecinos centrales y anteriores bien circunscritos,
cámara anterior sin células ni flare, iris claro con transiluminación, cristalino transparente e
hipoplasia foveal. Con los datos clínicos oculares y sistémicos se establecieron los diagnósticos
de AOC y DG y se citó a la familia para valoración clínica.

Se trata de una familia de nueve miembros, con padres de 74 y 64 años. No hay consanguinidad y
proceden de localidades diferentes. La DG está presente en la madre, tres hermanos y su hijo,
mientras que el AOC aparece en dos de los hermanos pero sin datos de DG.

\section{Caso 6}\label{anexo:texto_caso_6}

Mujer de 77 años de edad con único antecedente de hipertensión arterial. Presentaba un cuadro
clínico consistente en ataxia de la marcha, incontinencia de esfínteres y trastorno de funciones
superiores. Los estudios de neuroimagen mostraron dilatación del sistema ventricular, por lo que,
ante la sospecha de una HCA, se procedió a la colocación de un DLE. A las 30 horas de su colocación
comenzó, bruscamente, a presentar cefaleas y vómitos no asociados a ortostatismo. No hubo deterioro
del nivel de consciencia. En la bolsa de drenaje había 510 ml de LCR. Una TC craneal, mostró un
hematoma en el hemisferio cerebeloso izquierdo, de disposición transversal, con escaso efecto de
masa. Ante los hallazgos y la situación clínica de la paciente, se decidió la retirada del DLE y
tratamiento conservador. En la RM cerebral realizada previa al ingreso se había descartado
patología subyacente. A los cinco días la paciente fue dada de alta con recuperación hasta su
estado premórbido.

No se le implantó válvula de derivación al no objetivarse mejoría con la colocación de DLE.

\section{Caso 7}\label{anexo:texto_caso_7}

Paciente varón, de 40 años de edad, con antecedentes de salud. Acude a consulta de Urología por
presentar disuria y 3 episodios de uretrorragia. Al examen físico las mucosas son húmedas y
normocoloreadas, la auscultación cardio-respiratoria es normal, el abdomen es negativo, fosas
lumbares libres, el tacto rectal, los genitales externos y las regiones inguinales son normales. Se
le realiza ecografía urológica que es normal. La cistouretroscopía es normal, sólo notándose un
leve resalto al paso del cistoscopio por uretra anterior. Se le realiza radiografía de
uretrocistografía miccional donde se observa un defecto de llenado a nivel de la uretra anterior.

Por todo lo anterior decidimos realizar un estudio endoscópico bajo anestesia, detectándose una
tumoración polipoide, de color rojiza, bien delimitada, por lo que se decide resección endoscópica
de la misma. El resultado anatomopatológico de la pieza quirúrgica fue un papiloma invertido de
uretra anterior. Se da de alta por mejoría, evolucionando satisfactoriamente con desaparición de la
disuria y la uretrorragia.

\section{Caso 8}\label{anexo:texto_caso_8}

Paciente varón de 69 años, con antecedente de gastrectomía parcial hace 30 años por ulcus gástrico
péptico con reconstrucción de tipo Billroth II que ingresa por clínica de ocho días de evolución
consistente en dolor abdominal localizado en mesogastrio junto a un leve aumento de las cifras de
amilasa y lipasa séricas (amilasa 550 unidades interlacionales por litro, lipasa 5976 unidades
interlacionales por litro). La ecografía abdominal mostró la vesícula biliar distendida junto a una
dilatación de las vías biliares intra y extrahepáticas y una gran formación quística de paredes
finas y morfología tubular ocupando desde hipocondrio izquierdo hasta vacío derecho, lo que sugería
que se tratase del asa aferente dilatada y replecionada de líquido. La TAC y la RMN confirmaron
estos hallazgos. La endoscopia digestiva alta permitió visualizar la existencia de pliegues
engrosados en el muñón gástrico que obstruían el asa aferente, lo que sugería infiltración
neoplásica que fue confirmada tras el estudio histológico de las biopsias endoscópicas. El paciente
fue intervenido quirúrgicamente, con gastrectomía total y anastomosis esófago-yeyunal.

\section{Caso 9}\label{anexo:texto_caso_9}

Niña de 13 años de edad cuando se realiza el diagnóstico de SDRC en ambas extremidades inferiores
sin traumatismo previo.

Derivada a la Unidad del Dolor por el Servicio de Traumatología Pediátrica a los cuatro meses del
diagnóstico. Se administró tratamiento farmacológico más aplicación del parche de capsaicina al 8\%
con remisión completa del cuadro.

\section{Caso 10}\label{anexo:texto_caso_10}

Niña de tres años de edad diagnosticada hace un año de neurofibromatosis tipo 1 por el servicio de
pediatría. Se remite a oftalmología para valoración, en el contexto del cribado de otras patologías
que se lleva a cabo en estos pacientes. Por ello, también se solicita una resonancia magnética
nuclear (RMN) craneal.

En este momento la paciente carecía de sintomatología alguna y la exploración era la siguiente:
agudeza visual sin corrección en el ojo derecho 0,9 y en el izquierdo 1. Exploración
biomicroscópica del segmento anterior era normal, así como la funduscopia en ambos ojos. Presión
intraocular normal en los dos ojos. Movimientos oculares externos, sin hallazgos patológicos. No
proptosis ni exoftalmos. Pupilas isocóricas y normorreactivas.

En el estudio con neuroimagen con RMN se detectó la presencia de un glioma del nervio óptico del
ojo izquierdo.

Se decidió no tratar, pero seguir a la paciente con revisiones periódicas, cada tres meses, por
nuestra parte. En dicho seguimiento, progresivamente ha aparecido una leve proptosis y la visión de
la paciente ha disminuido a 0,9 en el ojo izquierdo.

\section{Ejemplos de informes generados}\label{anexo:ejemplos_informes_generados}

Esta sección reproduce, sin ninguna edición, tres informes clínicos reales generados por el
\englishText{LLM} del sistema (\texttt{Informe.txt}) a partir de los textos de referencia
anteriores, seleccionados por ilustrar directamente los patrones de error discutidos en la
Sección~\ref{section:resultados_e2e}: la propagación de un error de transcripción hasta un mapeo
terminológico incorrecto (\texttt{Test2}, \textit{batch}; Listado~\ref{lst:informe_test2_batch}),
la alucinación de un diagnóstico y de un tratamiento farmacológico inexistente en el dictado
original (\texttt{Test9}, \textit{batch}; Listado~\ref{lst:informe_test9_batch}), y un informe
aceptado bajo el criterio del \englishText{LLM-as-judge}, sin errores graves
(\texttt{Test7}, \textit{streaming}; Listado~\ref{lst:informe_test7_streaming}).

\begin{lstlisting}[
    caption={Informe clínico generado para \texttt{Test2} (\textit{batch}, Anexo~\ref{anexo:texto_caso_2}). El motor de normalización mapea la sigla mal transcrita ``CTG'' a ``Cardiotochogram reactivity'' y ``Positiva'' a ``Mantoux: positive'', dos conceptos SNOMED CT ajenos al caso.},
    label={lst:informe_test2_batch},
    basicstyle=\scriptsize\ttfamily,
    breaklines=true,
    frame=single,
    numbers=none
]
**INFORME CLÍNICO**

* **Motivo de Consulta / Enfermedad Actual:** Paciente varón de 29 años, con antecedentes de
amilatisis en la infancia y operado a los 6 años. Presenta hábito intestinal normal, analítica
normal con ligera elevación persistente en las transaminazas, y resultados positivos para DQ2
y negativos para DQ8.
* **Plan Diagnóstico y Terapéutico:**
 + Medicamentos: No se mencionan medicamentos específicos.
 + Exámenes solicitados: Biopsía duodenal (ya realizada).
* **Impresión Diagnóstica:** Celiacia, en estudio.

**ANEXO DE CODIFICACIÓN CLÍNICA (SNOMED CT)**

| Entidad | Concept ID | Término Preferido |
| --- | --- | --- |
| Fumador | 77176002 | Smoker (en) |
| Infancia | 255398004 | Childhood (en) |
| DQ2 | 117164003 | HLA-DQ2 antigen (en) |
| DQ8 | 260068001 | Human leukocyte antigen DQ8(3) (en) |
| CTG | 251667000 | Cardiotochogram reactivity (en) |
| Positiva | 268376005 | Mantoux: positive (en) |
| Atrofia moderada | 771307003 | Charcot-Marie-Tooth disease type 2B5 (en) |
| Vellosidades intestinales | 23230007 | Intestinal villus (en) |
| Lámina | 88446008 | Laminar (en) |

Nota: Se han excluido las entidades que no tienen concepto SNOMED coincidente o correspondencia CIE-10.
\end{lstlisting}

\begin{lstlisting}[
    caption={Informe clínico generado para \texttt{Test9} (\textit{batch}, Anexo~\ref{anexo:texto_caso_9}). El LLM expande la sigla truncada ``DSDRC'' como ``Distrofia de Duchenne de Becker'', sustituyendo el diagnóstico real (síndrome de dolor regional complejo), y alucina un plan terapéutico con seis fármacos ausentes del dictado original.},
    label={lst:informe_test9_batch},
    basicstyle=\scriptsize\ttfamily,
    breaklines=true,
    frame=single,
    numbers=none
]
**INFORME CLÍNICO**

* **Motivo de Consulta / Enfermedad Actual:** DSDRC (Distrofia de Duchenne de Becker) en ambas
extremidades inferiores sin traumatismo previo. La niña de 13 años ha sido tratada con remisión
completa del cuadro con aplicación del parche de capsaicina al 8% y tratamiento farmacológico.
* **Plan Diagnóstico y Terapéutico:** Se han prescrito los siguientes medicamentos:
 + Enalapril
 + Losartán (x2)
 + Omeprazol (x2)
 + Metformina
 + Atorvastatina
 + Capasaicina al 8% en parche

**ANEXO DE CODIFICACIÓN CLÍNICA (SNOMED CT)**

| Entidad | Concept ID | Término Preferido |
| --- | --- | --- |
| Traumatismo | 417746004 | Traumatic injury |
| Capsaicina | 400416009 | Capsaicin-containing product |

Nota: Se han eliminado los códigos de medicamentos duplicados y solo se ha incluido el término
preferido para cada entidad.
\end{lstlisting}

\begin{lstlisting}[
    caption={Informe clínico generado para \texttt{Test7} (\textit{streaming}, Anexo~\ref{anexo:texto_caso_7}), uno de los tres informes aceptados en esta modalidad bajo el criterio del LLM-as-judge (Sección~\ref{section:resultados_e2e}).},
    label={lst:informe_test7_streaming},
    basicstyle=\scriptsize\ttfamily,
    breaklines=true,
    frame=single,
    numbers=none
]
La información proporcionada parece ser un informe de caso médico que describe el tratamiento y
evolución de una paciente con un tumoración polipoide en la uretra anterior. A continuación, se
resumen los puntos clave del informe:

1. **Antecedentes**: La paciente presentó síntomas de disuria (dolor al orinar) y uretrorragia
(sangrado de la uretra).
2. **Exploración**: Se realizó un estudio endoscópico bajo anestesia, que detectó una tumoración
polipoide en la uretra anterior.
3. **Diagnóstico**: El diagnóstico fue un papiloma invertido en la uretra anterior.
4. **Tratamiento**: La paciente se sometió a una resección endoscópica de la tumoración.
5. **Resultados**: El resultado anatomopatológico de la pieza quirúrgica fue un papiloma
invertido, y la paciente evolucionó satisfactoriamente con desaparición de los síntomas.

Las entidades identificadas en el informe son:

1. **Dysuria** (disuria): dolor al orinar.
2. **Uretrorragia**: sangrado de la uretra.
3. **Papiloma invertido**: un tipo de tumor que puede ocurrir en la uretra.
4. **Resección endoscópica**: una técnica quirúrgica utilizada para eliminar el tumor.
5. **Anestesia**: medicamento utilizado para aliviar el dolor durante el procedimiento.

En cuanto a las correspondencias con CIE-10 (Clasificación Internacional de Enfermedades, 10na
revisión), se observa que algunas de las entidades identificadas no tienen una correspondencia
directa en la clasificación. Esto puede deberse a diferentes nombres o descripciones utilizados
en el informe en comparación con la clasificación CIE-10.

Es importante destacar que la información proporcionada es un resumen del informe y puede no ser
exhaustiva ni precisa al 100%.
\end{lstlisting}
